\chapter{Thoracoscopy and VATS}

\section*{What Is This Surgery?}
Thoracoscopy, specifically video-assisted thoracoscopic surgery (VATS), is a minimally invasive surgical technique employed for both diagnostic and therapeutic purposes within the thoracic cavity. This procedure is particularly significant for conditions affecting the pleura, lungs, mediastinum, diaphragm, and occasionally the pericardium. By utilizing small incisions and specialized instruments, VATS allows surgeons to visualize and operate on internal thoracic structures with the aid of a high-definition video camera. The clinical importance of VATS lies in its ability to provide a less traumatic alternative to traditional open thoracotomy, resulting in reduced postoperative pain, shorter hospital stays, and quicker recovery times for patients. This approach has transformed the management of various thoracic diseases, making it a cornerstone of modern thoracic surgery.

\section*{Alternative Names}
\begin{itemize}
  \item VATS
  \item Video-assisted thoracoscopic surgery
  \item Pleuroscopy
  \item Medical thoracoscopy
  \item Minimally invasive thoracic surgery (MITS)
  \item Thoracoscopic surgery
\end{itemize}

\section*{Relevant Surgical Anatomy}
The thoracic cavity houses critical structures, including the lungs, pleura, mediastinum, diaphragm, and major vessels. The lungs are divided into lobes, with the right lung having three lobes and the left lung having two. The pleura consists of two layers: the visceral pleura covering the lungs and the parietal pleura lining the thoracic cavity. The arterial supply to the lungs is primarily through the pulmonary arteries, while venous drainage occurs via the pulmonary veins. The bronchial arteries supply oxygenated blood to the lung tissue.

Key anatomical landmarks include:

\begin{itemize}
  \item \textbf{McBurney's Point:} Although primarily associated with appendicitis, this landmark is crucial for understanding abdominal anatomy in thoracic procedures.
  \item \textbf{Calot's Triangle:} Important in thoracic surgeries involving the mediastinum, this triangle is formed by the cystic duct, common hepatic duct, and the liver.
  \item \textbf{Phrenic Nerve:} This nerve innervates the diaphragm and is critical to monitor during thoracoscopic procedures to avoid diaphragmatic paralysis.
  \item \textbf{Intercostal Nerves and Vessels:} These structures run along the inferior border of each rib and are vital to avoid during port placement.
\end{itemize}

Understanding these anatomical relationships is essential for safe and effective thoracoscopic surgery, as they guide the surgeon in avoiding complications and ensuring optimal outcomes.

\section*{Surgical Principle \& Rationale}
The primary principle of VATS is to access the thoracic cavity through multiple small incisions, typically ranging from 1 to 5 centimeters, rather than performing a large thoracotomy. This technique minimizes trauma to the chest wall and reduces postoperative pain. The use of a thoracoscope allows for a magnified view of the surgical field, enabling precise dissection and manipulation of tissues. The rationale behind VATS is to achieve effective diagnosis and treatment while preserving lung function and expediting recovery. The technical goals include complete resection of diseased tissue, accurate staging of malignancies, and effective palliation of symptoms, all while minimizing the physiological stress associated with traditional open surgery.

\section*{History \& Surgical Evolution}
The history of thoracoscopy dates back to 1910 when Swedish internist Hans Christian Jacobaeus performed the first documented pleuroscopy using a cystoscope to treat pleural adhesions in tuberculosis patients. For many years, thoracoscopy was primarily a diagnostic tool utilized by pulmonologists. The introduction of fiber optics in the 1970s enhanced visualization, but it was the advent of video technology in the late 1980s and early 1990s that revolutionized thoracic surgery. Pioneers like Giancarlo Roviaro were instrumental in applying video-assisted techniques to complex lung resections, such as lobectomies. Over the years, VATS has evolved from a three-port technique to more refined approaches, including two-port and uniportal methods, establishing itself as the standard for many thoracic surgical procedures.

\section*{Clinical Indications}
VATS is indicated for a variety of clinical scenarios, including:

\begin{itemize}
  \item Diagnosis of undiagnosed pleural effusions, particularly those suspicious for malignancy or infection.
  \item Biopsy of indeterminate lung nodules or masses to establish a definitive diagnosis.
  \item Staging and definitive resection of early-stage lung cancer, including lobectomy and segmentectomy.
  \item Drainage and decortication for empyema or complicated parapneumonic effusions.
  \item Treatment of recurrent or spontaneous pneumothorax, often involving pleurodesis or bullectomy.
  \item Biopsy or resection of mediastinal masses, such as thymomas or neurogenic tumors.
  \item Creation of a pericardial window for recurrent pericardial effusions.
  \item Sympathectomy for hyperhidrosis.
  \item Management of benign esophageal conditions, such as achalasia.
\end{itemize}

These indications reflect the versatility of VATS in addressing both diagnostic and therapeutic needs in thoracic surgery.

\section*{Clinical Positioning}
VATS has transitioned from a secondary or adjunctive procedure to a primary surgical modality for many thoracic conditions. For early-stage lung cancer, VATS lobectomy is now considered the preferred approach, offering oncologic outcomes comparable to open thoracotomy with reduced morbidity. It serves as a first-line option for conditions such as recurrent pneumothorax and undiagnosed pleural effusions. In cases where initial non-surgical treatments have failed, VATS may be employed as a secondary intervention. Additionally, VATS can be utilized as an initial diagnostic step, with the possibility of conversion to an open thoracotomy if significant adhesions or complications arise.

\section*{Surgical Technique \& Procedure}
VATS is performed under general anesthesia, often requiring single-lung ventilation to collapse the operative lung for optimal visualization. The patient is typically positioned in a lateral decubitus position, with the arm on the operative side elevated.

The procedure generally involves the following steps:

\begin{itemize}
  \item \textbf{Anesthesia and Positioning:} General anesthesia is induced, and a double-lumen endotracheal tube is placed for selective lung ventilation. The patient is positioned appropriately.
  \item \textbf{Port Placement:} One to three small incisions are made in the chest wall. The first port, usually in the mid-axillary line, allows for the insertion of the thoracoscope for diagnostic exploration.
  \item \textbf{Diagnostic Survey:} A thorough inspection of the lung, pleura, diaphragm, and mediastinum is performed to confirm pathology and plan the surgical strategy.
  \item \textbf{Working Ports Creation:} Additional ports are created under direct vision to introduce surgical instruments for dissection and manipulation.
  \item \textbf{Operative Steps:} Depending on the procedure, this may involve dissection of lung tissue, division of blood vessels and bronchi, lymph node dissection, or drainage of pleural fluid.
  \item \textbf{Specimen Retrieval:} Resected tissue is placed in an endoscopic retrieval bag and removed through one of the port sites.
  \item \textbf{Hemostasis and Irrigation:} The surgical field is checked for bleeding, and the chest cavity is irrigated with saline.
  \item \textbf{Chest Tube Placement:} A chest tube is inserted to drain air and fluid from the pleural space.
  \item \textbf{Closure:} The port site incisions are closed in layers with sutures, and sterile dressings are applied.
\end{itemize}

This step-by-step approach ensures a thorough and effective surgical intervention while minimizing trauma to surrounding tissues.

\section*{Preoperative Workup \& Preparation}
Thorough preoperative preparation is essential for optimizing outcomes in VATS. This typically includes:

\begin{itemize}
  \item \textbf{Medical History and Physical Examination:} Comprehensive assessment of the patient's overall health and comorbidities.
  \item \textbf{Laboratory Tests:} Routine blood work, including complete blood count (CBC), electrolyte panel, renal and liver function tests, and coagulation studies.
  \item \textbf{Imaging Studies:} A recent chest X-ray and high-resolution computed tomography (CT) scan of the chest are mandatory.
  \item \textbf{Pulmonary Function Tests (PFTs):} Essential for assessing lung capacity and ability to tolerate the procedure.
  \item \textbf{Cardiac Evaluation:} An electrocardiogram (ECG) is standard, with further assessment as needed.
  \item \textbf{Anesthesia Consultation:} To evaluate the patient's fitness for general anesthesia.
  \item \textbf{Medication Management:} Instructions on which medications to stop or continue before surgery.
  \item \textbf{NPO Status:} Patients are typically instructed to be NPO for 6-8 hours prior to surgery.
  \item \textbf{Preoperative Antibiotics:} Administered intravenously shortly before the incision to prevent infection.
  \item \textbf{Informed Consent:} Detailed discussion with the patient regarding the procedure, risks, benefits, and alternatives.
\end{itemize}

This comprehensive preparation helps ensure patient safety and optimal surgical outcomes.

\section*{Contraindications \& Precautions}
While VATS is widely applicable, certain conditions may contraindicate its use or necessitate caution.

\textbf{Absolute Contraindications:}
\begin{itemize}
  \item Inability to tolerate single-lung ventilation due to severe lung disease.
  \item Severe, uncorrectable coagulopathy.
  \item Hemodynamic instability that precludes general anesthesia.
  \item Extensive pleural adhesions preventing safe entry into the pleural space.
\end{itemize}

\textbf{Relative Contraindications and Precautions:}
\begin{itemize}
  \item Prior extensive thoracic surgery or radiation therapy.
  \item Large tumor size or invasion of critical mediastinal structures.
  \item Severe underlying lung disease that increases the risk of complications.
  \item Morbid obesity affecting positioning and instrument manipulation.
  \item Lack of surgeon experience with complex VATS procedures.
  \item Inability to localize small, non-palpable lung nodules preoperatively.
\end{itemize}

Awareness of these contraindications is crucial for ensuring patient safety and optimizing surgical outcomes.

\section*{Complications \& Risks}
As with any surgical procedure, VATS carries potential risks, which can be categorized into general surgical risks and procedure-specific complications.

\begin{itemize}
  \item \textbf{General Surgical Risks:}
    \begin{itemize}
      \item \textbf{Bleeding (Hemorrhage):} Ranging from minor oozing to significant blood loss requiring transfusion or re-operation.
      \item \textbf{Infection:} Including surgical site infection, pneumonia, or empyema.
      \item \textbf{Anesthesia Complications:} Adverse reactions to anesthetic agents, cardiac events, or respiratory complications.
      \item \textbf{Pain:} Postoperative pain, though generally less severe than with open thoracotomy.
    \end{itemize}
  \item \textbf{Procedure-Specific Risks:}
    \begin{itemize}
      \item \textbf{Prolonged Air Leak:} Air leaking from the lung into the chest tube for more than 5-7 days post-surgery.
      \item \textbf{Pneumothorax:} Persistent or recurrent lung collapse.
      \item \textbf{Injury to Lung Parenchyma:} Accidental damage to healthy lung tissue.
      \item \textbf{Injury to Major Vessels or Nerves:} Damage leading to complications such as diaphragmatic paralysis or vocal cord dysfunction.
      \item \textbf{Chylothorax:} Leakage of lymphatic fluid into the pleural space.
      \item \textbf{Bronchopleural Fistula:} Abnormal connection between a bronchus and the pleural space.
      \item \textbf{Arrhythmias:} Particularly atrial fibrillation post-thoracic surgery.
      \item \textbf{Post-thoracotomy Pain Syndrome:} Chronic neuropathic pain along incision sites.
      \item \textbf{Tumor Seeding at Port Sites:} Rare complication in oncologic resections.
      \item \textbf{Incomplete Resection or Diagnosis:} Potentially requiring further intervention.
    \end{itemize}
\end{itemize}

Understanding these risks is essential for informed consent and patient management.

\section*{Postoperative Recovery Timeline}
Following VATS, patients are typically transferred to the Post-Anesthesia Care Unit (PACU) for close monitoring of vital signs, pain levels, and chest tube function. Immediate postoperative care focuses on pain control, respiratory support, and early mobilization.

Within the first 24-48 hours, patients are encouraged to ambulate to prevent complications like pneumonia and deep vein thrombosis. Chest tube management is critical; the tube drains air and fluid, and its output and presence of air leak are closely monitored. The chest tube is typically removed when drainage is minimal and there is no persistent air leak. Pain management strategies often include multimodal approaches to facilitate comfort and respiratory effort. Most patients are discharged within 1 to 5 days, depending on the complexity of the procedure and their recovery trajectory. For the first 1-2 weeks at home, patients are advised to gradually increase their activity levels, perform breathing exercises, and care for their wound sites. Heavy lifting and strenuous activities are generally restricted for several weeks. Long-term recovery typically sees a full return to normal activities within 2-4 weeks, significantly faster than recovery from an open thoracotomy.

\section*{Surgical Findings \& Pathology}
During VATS, surgeons document various intraoperative findings, including the condition of the pleura, lung tissue, and any masses or lesions. The pathology report on resected tissues provides critical information regarding the nature of the disease, such as the presence of malignancy, type of infection, or benign conditions. This report is essential for guiding further treatment decisions and follow-up care.

\section*{Expected Postoperative Course}
A normal postoperative course following VATS is characterized by the resolution of pain, improvement in respiratory function, and gradual return to baseline activities. Patients typically experience a decrease in symptoms such as shortness of breath or chest pain, with most reporting significant relief within days of the procedure. The timeline for recovery varies based on the complexity of the surgery and the patient's overall health, but many patients can resume normal activities within 2-4 weeks.

\section*{Postoperative Care \& Follow-up}
Postoperative follow-up is crucial for monitoring recovery and addressing any concerns. Key aspects of postoperative care include:

\begin{itemize}
  \item \textbf{Postoperative Clinic Visits:} Scheduled 1-2 weeks after discharge for wound assessment and review of pathology results.
  \item \textbf{Suture/Staple Removal:} If non-absorbable sutures or staples were used, they are removed during the initial follow-up visit.
  \item \textbf{Postoperative Imaging:} A chest X-ray is often performed at the first follow-up to confirm lung re-expansion.
  \item \textbf{Pulmonary Rehabilitation/Physical Therapy:} May be recommended for optimizing lung function and physical endurance.
  \item \textbf{Activity Resumption:} Gradual return to normal activities is advised, with specific instructions on avoiding strenuous exercise.
  \item \textbf{Oncology Follow-up:} For cancer patients, ongoing care with an oncologist is essential for adjuvant therapy planning.
  \item \textbf{Warning Signs:} Patients are educated on warning signs that warrant immediate medical attention, such as fever, increasing chest pain, or significant changes in chest tube site appearance.
\end{itemize}

This structured follow-up ensures that any complications are promptly addressed and that patients receive appropriate ongoing care.

\section*{Epidemiology}
VATS has seen a dramatic increase in utilization since its introduction, becoming the predominant surgical approach for many thoracic pathologies. For lung cancer resections, particularly lobectomy and segmentectomy, VATS now accounts for an estimated 60-80\% of procedures in many high-volume centers. Spontaneous pneumothorax frequently affects young, tall males, with VATS pleurodesis or bullectomy offering definitive treatment. Malignant pleural effusions, often seen in older patients with advanced cancers, are also frequently managed with VATS. Morbidity rates associated with VATS are generally lower than those for open thoracotomy, typically ranging from 10-30\% depending on the complexity of the procedure and patient comorbidities. Mortality rates for VATS lobectomy are reported to be between 1-3\%, underscoring the safety profile of this minimally invasive approach.

\section*{Outcomes \& Prognosis}
VATS procedures are associated with excellent outcomes and a favorable prognosis for a wide range of thoracic conditions. Patients undergoing VATS typically experience significantly reduced postoperative pain, shorter hospital stays, and faster recovery compared to traditional open thoracotomy. For early-stage lung cancer, numerous studies have demonstrated that oncologic outcomes, including disease-free survival and overall survival, are comparable to or even superior to open surgery. Recurrence rates for conditions like spontaneous pneumothorax are substantially reduced after VATS pleurodesis or bullectomy. Prognostic factors are largely dependent on the underlying disease; for cancer, tumor stage, nodal involvement, and completeness of resection are paramount determinants of long-term survival. Long-term functional outcomes, particularly pulmonary function, are often better preserved after VATS due to less disruption of the chest wall musculature and diaphragm.

\section*{References}
\begin{enumerate}
  \item MedlinePlus. Video-assisted thoracoscopic surgery. U.S. National Library of Medicine; 2025. Available from: \url{https://medlineplus.gov/ency/article/007357.htm}
  
  \item Sabiston Textbook of Surgery: The Biological Basis of Modern Surgical Practice. 22nd ed. Elsevier; 2022.
  
  \item Shields' General Thoracic Surgery. 8th ed. Wolters Kluwer; 2019.
  
  \item American College of Chest Physicians (ACCP) Guidelines for Management of Lung Cancer. Chest; 2021.
\end{enumerate}

\clearpage